% Section 5: Discussion
\section{Discussion}

%% ---- Slide 1: Key Findings Summary ----
\begin{frame}{Key Findings Summary}
  \begin{block}{What the Results Mean}
    \begin{itemize}
      \item \highlight{No universally best method} exists across 11 pathogens
        \begin{itemize}
          \item LOPO cross-validation: \highlight{0/11 match rate}
          \item Friedman test: $p = 0.990$ --- no method reliably superior
        \end{itemize}
      \item Scoring $\times$ pathogen interaction dominates performance variance
        \begin{itemize}
          \item $\omega^2 = 0.296$ (large effect), confirmed at 6-category level ($\omega^2 = 0.103$)
        \end{itemize}
      \item Multi-source integration recovers positions missed by single-scoring
        \begin{itemize}
          \item H3N2 EqualWeight: 4.012$\times$ escape enrichment ($p = 0.0023$)
          \item 4 core features capture 98\% of full-ensemble performance
        \end{itemize}
      \item Dual-benchmark design (synthetic + real) is itself a methodological contribution
        \begin{itemize}
          \item Real GISAID: enrichment 6.9--8.0$\times$ above random despite F1 drop
        \end{itemize}
    \end{itemize}
  \end{block}
\end{frame}

%% ---- Slide 2: Comparison with Existing Approaches ----
\begin{frame}{Comparison with Existing Approaches}
  \small
  \begin{tabular}{p{2.2cm}p{4.5cm}p{5.5cm}}
    \toprule
    \textbf{Prior Work} & \textbf{Their Finding} & \textbf{MutBench Extension} \\
    \midrule
    EVEREST & VEP performance varies across 45 viral DMS datasets &
      Convergent finding despite different tasks (fitness prediction vs.\ hotspot detection) \\[3pt]
    Maher et al. & Combining features outperforms single features (SARS-CoV-2) &
      Extended to \highlight{11 pathogens}; multi-source benefit generalizes \\[3pt]
    Hie et al. & PLM semantic similarity captures escape potential &
      Confirmed for EV-A71 (MCC=0.260), \highlight{not universal} \\[3pt]
    Bailey et al. & Cancer driver tools disagree on gene lists &
      Same pattern in virology: disagreement structured by pathogen biology \\[3pt]
    ProteinGym & VEP varies across 217 DMS assays &
      \highlight{Information type} (not just model) is the dominant factor \\
    \bottomrule
  \end{tabular}

  \vspace{0.5em}
  \begin{block}{Unifying Principle}
    Pathogen-dependence of optimal methods is a \highlight{general principle} ---
    MutBench quantifies it via variance decomposition ($\omega^2 = 0.296$)
  \end{block}
\end{frame}

%% ---- Slide 3: Implications for Mutation Analysis / Public Health ----
\begin{frame}{Implications for Public Health Surveillance}
  \begin{columns}[T]
    \column{0.48\textwidth}
    \begin{block}{Vaccine Escape Validation}
      \begin{itemize}
        \item H3N2: \highlight{9.36$\times$} enrichment ($p < 0.0001$)
        \item HIV-1: 7.19$\times$ enrichment ($p < 0.0001$)
        \item SARS-CoV-2: 7.63$\times$ ($p = 0.026$)
        \item Layer~A rankings concordant with escape enrichment --- \highlight{not circular}
      \end{itemize}
    \end{block}

    \column{0.48\textwidth}
    \begin{block}{Practical Workflow for New Pathogens}
      \begin{enumerate}
        \item Compute 20 scoring types from MSA
        \item Use family-level recommendation if available
        \item Otherwise: EqualWeight ensemble (4 core features)
        \item Detect with KDE/Wavelet defaults
      \end{enumerate}
      \vspace{0.3em}
      {\footnotesize Runtime: $\sim$15--20 min per pathogen (single GPU)}
    \end{block}
  \end{columns}

  \vspace{0.5em}
  \centering
  \begin{tabular}{lll}
    \toprule
    \textbf{Virus Family} & \textbf{Recommended Scoring} & \textbf{Rationale} \\
    \midrule
    Orthomyxoviridae & Phylogenetic (FUBAR) & Diversifying selection \\
    Coronaviridae    & AI-based (PLM)       & Shallow phylogeny \\
    Retroviridae     & MSA-derived (entropy) & High diversity \\
    Flaviviridae     & Frequency-based       & Hypervariable regions \\
    \bottomrule
  \end{tabular}
\end{frame}

%% ---- Slide 4: Limitations and Threats to Validity ----
\begin{frame}{Limitations and Threats to Validity}
  \begin{columns}[T]
    \column{0.48\textwidth}
    \begin{block}{Methodological Limitations}
      \begin{itemize}
        \item \highlight{H-score saturation}: all 1,273 Spike positions nonzero in 2024--2025 (SD = 0.098)
        \item \highlight{Phylogenetic non-independence}: frequency cannot distinguish selection from founder effects (D614G)
        \item Bootstrap CIs wide at region level ($n=7$), but pairwise comparisons significant ($p < 0.001$)
        \item ESM-2 training data overlap with benchmark proteins
      \end{itemize}
    \end{block}

    \column{0.48\textwidth}
    \begin{block}{Data and Scope Limitations}
      \begin{itemize}
        \item \highlight{11 RNA viruses} only (DNA viruses need different scoring)
        \item 3 family pairs not phylogenetically independent (but within-family best methods differ)
        \item GISAID geographic bias: $>$80\% from high-income countries
        \item Ground truth heterogeneity across pathogens (convergent evolution vs.\ HVR vs.\ epochal)
        \item Low recall for short proteins: EV-A71 (0.185), Rabies (0.308)
      \end{itemize}
    \end{block}
  \end{columns}
\end{frame}

%% ---- Slide 5: Broader Impact and Theoretical Framing ----
\begin{frame}{Broader Impact}
  \begin{block}{Theoretical Contribution: Rice's Algorithm Selection Problem}
    \begin{itemize}
      \item Problem space $\mathcal{P}$: pathogens with genomic profiles
      \item Algorithm space $\mathcal{A}$: 780 scoring--detection combinations
      \item Mapping $f: \mathcal{P} \to \mathcal{A}$ is \highlight{non-trivial} ($\omega^2 = 0.296$)
      \item Empirically grounded No Free Lunch --- stronger than abstract theorem
    \end{itemize}
  \end{block}

  \begin{block}{Three Problems Addressed (from Chapter 1)}
    \begin{enumerate}
      \item \highlight{No standardized evaluation} $\to$ 3-layer ground truth + hotspot-score metric
      \item \highlight{No cross-pathogen evidence} $\to$ 4 independent lines of statistical evidence
      \item \highlight{Unknown pathogen-dependence} $\to$ $\omega^2 = 0.296$ interaction + LOPO 0/11 failure
    \end{enumerate}
  \end{block}

  \begin{block}{Key Takeaway}
    MutBench enables \highlight{pathogen-adaptive method selection} ---
    reducing experimental search space rather than achieving perfect prediction
  \end{block}
\end{frame}
